%All acronyms must be written in this file.


\frontmatter % Use roman page numbering style (i, ii, iii, iv...) for the pre-content pages

\pagestyle{plain} % Default to the plain heading style until the thesis style is called for the body content

%----------------------------------------------------------------------------------------
%	COVER PAGE
%----------------------------------------------------------------------------------------

\makecoverpage

%----------------------------------------------------------------------------------------
%	TITLE PAGE
%----------------------------------------------------------------------------------------

\maketitlepage


%----------------------------------------------------------------------------------------
%	STATEMENT of INTEGRITY
%----------------------------------------------------------------------------------------
\integritystatement

%----------------------------------------------------------------------------------------
%	ABSTRACT PAGE
%----------------------------------------------------------------------------------------

\begin{abstract}
% Resumo alargado em português (~1000 palavras) — obrigatório para dissertações escritas em inglês.

As organizações de saúde digital realizam frequentemente tarefas manuais e repetitivas no âmbito dos seus processos internos de investigação e desenvolvimento (I\&D), tais como a documentação de projetos, a organização de espaços de trabalho e a geração de secções recorrentes em relatórios. Estas tarefas consomem tempo significativo às equipas, limitando a capacidade de se concentrarem em trabalho de maior valor acrescentado.

Os avanços recentes em modelos de linguagem de grande escala (LLM — \textit{Large Language Models}) e nas arquiteturas de agentes inteligentes abrem novas possibilidades para automatizar este tipo de fluxos de trabalho. Em particular, a combinação de agentes LLM com ferramentas externas e com Geração Aumentada por Recuperação (RAG — \textit{Retrieval-Augmented Generation}) permite criar sistemas capazes de aceder a informação contextual, executar ações sobre sistemas externos e gerar conteúdo fundamentado em documentação existente.

No entanto, o domínio da saúde impõe restrições específicas que tornam a automação direta inviável sem salvaguardas adequadas. Entre os principais riscos destacam-se as alucinações dos modelos, a execução de ações irreversíveis sem validação humana, e as obrigações de rastreabilidade e privacidade decorrentes do Regulamento Geral sobre a Proteção de Dados (RGPD) e do Regulamento Europeu de Inteligência Artificial (AI Act). A literatura existente sobre autonomia humano-máquina (Parasuraman et al., 2000; Shneiderman, 2020) fornece modelos conceptuais úteis, mas é escassa em orientações práticas sobre como operacionalizar estes princípios em contextos reais de I\&D.

Esta dissertação responde à seguinte pergunta central: \textit{como conceber agentes LLM que automatizem processos internos de I\&D em saúde digital, mantendo supervisão humana e rastreabilidade, com ganhos mensuráveis de eficiência sem perda de qualidade nem de confiança por parte da equipa?} Como questão secundária, avalia-se o impacto na eficiência, na qualidade dos outputs e na aceitação dos utilizadores.

O trabalho é desenvolvido no contexto da VirtualCare, empresa portuguesa de saúde digital, onde são identificados dois a três processos de I\&D representativos. Para cada processo selecionado, é definido um modelo de autonomia que classifica cada ação possível do agente em três categorias: ações autónomas (executadas sem intervenção humana), ações sujeitas a aprovação (o agente propõe, o humano valida) e ações proibidas (nunca executadas de forma autónoma). Este modelo operacionaliza o conceito de \textit{Human-in-the-Loop} (HITL), garantindo que a supervisão humana ocorre precisamente nos pontos críticos de cada fluxo.

A metodologia adotada é a \textit{Design Science Research} (Hevner et al., 2004; Peffers et al., 2007), que enquadra o trabalho como a conceção e avaliação de um artefacto (o protótipo dos agentes) em resposta a um problema prático. O processo inclui: levantamento por entrevistas junto da equipa da VirtualCare para mapeamento dos processos; desenho do modelo de autonomia e da arquitetura dos agentes; implementação de um protótipo funcional para dois processos; e ensaios comparativos com a equipa.

O protótipo é avaliado segundo critérios quantitativos e qualitativos. Em termos de eficiência, é estabelecida uma meta de redução mínima de 30\% do tempo face ao processo manual. A qualidade dos outputs é medida por correção e completude (meta: $\geq$85\%) e por capacidade de deteção de erros injetados (meta: $\geq$80\%). A aceitação dos utilizadores é avaliada através do \textit{System Usability Scale} (SUS, Brooke, 1996), com limiar de 68, e de uma escala de utilidade percebida com meta $\geq$4/5.

São comparadas pelo menos três configurações: um fluxo fixo baseado em prompts sem autonomia do agente, um agente com ferramentas e RAG com nível de autonomia máxima, e variantes intermédias com diferentes níveis de supervisão humana. O baseline é o processo manual existente na VirtualCare.

A rastreabilidade é assegurada através de registos de auditoria que cobrem 100\% das ações do agente, sem exceções. Nenhuma ação irreversível é executada sem aprovação prévia de um utilizador humano. Os dados utilizados nos ensaios são anonimizados ou sintéticos, de acordo com as restrições de confidencialidade aplicáveis.

O trabalho contribui com orientações práticas e replicáveis para a integração de agentes LLM em fluxos de trabalho organizacionais sensíveis, com supervisão humana adequada. Estas orientações são concebidas de forma a serem generalizáveis a outros contextos de saúde digital com características semelhantes.

% TODO: expandir para ~1000 palavras com detalhes da implementação e resultados (a preencher no 2.º semestre)

\end{abstract}

\begin{abstractotherlanguage}

Organisations in digital health routinely perform time-consuming manual tasks as part of their research and development (R\&D) processes — including project documentation, workspace organisation, and the generation of recurring report sections. Large Language Model (LLM) agents equipped with tools and Retrieval-Augmented Generation (RAG) offer a promising approach for automating these workflows. However, the healthcare domain introduces critical constraints: model hallucinations, the risk of irreversible actions, strict traceability requirements, and privacy obligations under the GDPR and the EU AI Act.

This dissertation proposes a framework for designing LLM-based agents that automate selected internal R\&D processes while preserving Human-in-the-Loop (HITL) supervision at critical decision points. Following a Design Science Research methodology, the proposed artefact defines an autonomy model that distinguishes autonomous actions, approval-required actions, and prohibited actions. The study is grounded in a real-world context at VirtualCare, a Portuguese digital health company, where two to three representative processes are identified, mapped, and targeted for automation.

The evaluation covers efficiency (target: $\geq$30\% time reduction over the manual baseline), output quality (target: $\geq$85\% correctness and completeness), error detection (target: $\geq$80\% of injected errors detected), and user acceptance (System Usability Scale $\geq$68, perceived utility $\geq$4/5). The study concludes with replicable guidelines for integrating LLM agents into sensitive organisational workflows with appropriate human oversight.

\end{abstractotherlanguage}

%----------------------------------------------------------------------------------------
%	ACKNOWLEDGEMENTS (optional)
%----------------------------------------------------------------------------------------

\begin{acknowledgements}

% TODO: a preencher.

\end{acknowledgements}

%----------------------------------------------------------------------------------------
%	LIST OF CONTENTS/FIGURES/TABLES PAGES
%----------------------------------------------------------------------------------------

\tableofcontents % Prints the main table of contents

\listoffigures % Prints the list of figures

\listoftables % Prints the list of tables

\iflanguage{portuguese}{
\renewcommand{\listalgorithmname}{Lista de Algoritmos}
}
\listofalgorithms % Prints the list of algorithms
\addchaptertocentry{\listalgorithmname}


\renewcommand{\lstlistlistingname}{List of Source Code}
\iflanguage{portuguese}{
\renewcommand{\lstlistlistingname}{Lista de C\'odigo}
}
\lstlistoflistings % Prints the list of listings (programming language source code)
\addchaptertocentry{\lstlistlistingname}


%----------------------------------------------------------------------------------------
%	ACRONYMS
%----------------------------------------------------------------------------------------

\newcommand{\listacronymname}{List of Acronyms}
\iflanguage{portuguese}{
\renewcommand{\listacronymname}{Lista de Acr\'onimos}
}

%Use GLS
\glsresetall
\printglossary[title=\listacronymname,type=\acronymtype,style=long]

%----------------------------------------------------------------------------------------
%	DONE
%----------------------------------------------------------------------------------------

\mainmatter % Begin numeric (1,2,3...) page numbering
\pagestyle{thesis} % Return the page headers back to the "thesis" style
